\documentclass{amsart}
% For dealing with references we use the comment environment
\usepackage{verbatim}
\newenvironment{reference}{\comment}{\endcomment}
%\newenvironment{reference}{}{}
\newenvironment{slogan}{\comment}{\endcomment}
\newenvironment{history}{\comment}{\endcomment}

% For commutative diagrams we use Xy-pic
\usepackage[all]{xy}

% We use 2cell for 2-commutative diagrams.
\xyoption{2cell}
\UseAllTwocells

% We use multicol for the list of chapters between chapters
\usepackage{multicol}

% This is generally recommended for better output
\usepackage{lmodern}
\usepackage[T1]{fontenc}

% For cross-file-references

% Package for hypertext links:
\usepackage{hyperref}

% For any local file, say "hello.tex" you want to link to please

% Theorem environments.
%
\theoremstyle{plain}
\newtheorem{theorem}[subsection]{Theorem}
\newtheorem{proposition}[subsection]{Proposition}
\newtheorem{lemma}[subsection]{Lemma}

\theoremstyle{definition}
\newtheorem{definition}[subsection]{Definition}
\newtheorem{example}[subsection]{Example}
\newtheorem{exercise}[subsection]{Exercise}
\newtheorem{situation}[subsection]{Situation}

\theoremstyle{remark}
\newtheorem{remark}[subsection]{Remark}
\newtheorem{remarks}[subsection]{Remarks}

\numberwithin{equation}{subsection}

% Macros
%
\def\lim{\mathop{\mathrm{lim}}\nolimits}
\def\colim{\mathop{\mathrm{colim}}\nolimits}
\def\Spec{\mathop{\mathrm{Spec}}}
\def\Hom{\mathop{\mathrm{Hom}}\nolimits}
\def\Ext{\mathop{\mathrm{Ext}}\nolimits}
\def\SheafHom{\mathop{\mathcal{H}\!\mathit{om}}\nolimits}
\def\SheafExt{\mathop{\mathcal{E}\!\mathit{xt}}\nolimits}
\def\Sch{\mathit{Sch}}
\def\Mor{\mathop{\mathrm{Mor}}\nolimits}
\def\Ob{\mathop{\mathrm{Ob}}\nolimits}
\def\Sh{\mathop{\mathit{Sh}}\nolimits}
\def\NL{\mathop{N\!L}\nolimits}
\def\CH{\mathop{\mathrm{CH}}\nolimits}
\def\proetale{{pro\text{-}\acute{e}tale}}
\def\etale{{\acute{e}tale}}
\def\QCoh{\mathit{QCoh}}
\def\Ker{\mathop{\mathrm{Ker}}}
\def\Im{\mathop{\mathrm{Im}}}
\def\Coker{\mathop{\mathrm{Coker}}}
\def\Coim{\mathop{\mathrm{Coim}}}

% Boxtimes
%
\DeclareMathSymbol{\boxtimes}{\mathbin}{AMSa}{"02}

%
% Macros for moduli stacks/spaces
%
\def\QCohstack{\mathcal{QC}\!\mathit{oh}}
\def\Cohstack{\mathcal{C}\!\mathit{oh}}
\def\Spacesstack{\mathcal{S}\!\mathit{paces}}
\def\Quotfunctor{\mathrm{Quot}}
\def\Hilbfunctor{\mathrm{Hilb}}
\def\Curvesstack{\mathcal{C}\!\mathit{urves}}
\def\Polarizedstack{\mathcal{P}\!\mathit{olarized}}
\def\Complexesstack{\mathcal{C}\!\mathit{omplexes}}
% \Pic is the operator that assigns to X its picard group, usage \Pic(X)
% \Picardstack_{X/B} denotes the Picard stack of X over B
% \Picardfunctor_{X/B} denotes the Picard functor of X over B
\def\Pic{\mathop{\mathrm{Pic}}\nolimits}
\def\Picardstack{\mathcal{P}\!\mathit{ic}}
\def\Picardfunctor{\mathrm{Pic}}
\def\Deformationcategory{\mathcal{D}\!\mathit{ef}}

\setlength{\emergencystretch}{3em}
\begin{document}
\title[Stacks Project, Tag 055Q]{472 --- The fibre component containing a section is locally a neighbourhood at a geometrically reduced fibre in flat finitely presented families}
\maketitle
\noindent Copyright (C) 2005--2025 Johan de Jong. Stacks Project authors. Source: \href{https://stacks.math.columbia.edu/tag/055Q}{Tag 055Q}.

\noindent Editorial difficulty note (not author text). Difficulty H2: Noetherian approximation must preserve the section, flatness and geometric reducedness before a constructible-neighbourhood criterion can be used. The proof realizes every generalization by a discrete valuation ring and controls connected components there, showing that the selected fibre component is stable under the necessary generalizations..

\noindent Editorial provenance note (not author text). History: excerpt from revision \texttt{a0444\allowbreak 6e57e\allowbreak c1fbc\allowbreak 252a8\allowbreak 71afc\allowbreak ec775\allowbreak 2fb28\allowbreak 07b14}, more-morphisms.tex, lines 8309--8384. Reference presentation was adapted; original-author context and core support blocks were retained or added. The mathematical wording is unchanged. Licensed under GNU FDL 1.2 or later; no invariant sections or cover texts. See ../licenses/COPYING. Context blocks reproduce author definitions and supporting results; they are not additional counted problems.

\begin{situation}
\label{situation-connected-along-section}
Here $f : X \to Y$ be a morphism of schemes, and
$s : Y \to X$ is a section of $f$.
For every $y \in Y$ we denote $X^0_y$ the connected component of $X_y$
containing $s(y)$. Finally, we set $X^0 = \bigcup_{y \in Y} X^0_y$.
\end{situation}

\begin{lemma}
\label{lemma-connected-along-section-open-neighbourhood}
Let $f : X \to Y$, $s : Y \to X$ be as in
Situation \ref{situation-connected-along-section}.
Let $y \in Y$ be a point.
Assume
\begin{enumerate}
\item $f$ is of finite presentation and flat, and
\item the fibre $X_y$ is geometrically reduced.
\end{enumerate}
Then $X^0$ is a neighbourhood of $X^0_y$ in $X$.
\end{lemma}

\begin{proof}
We may replace $Y$ with an affine open neighbourhood of $y$.
Write $Y = \Spec(A)$ and $A = \colim A_i$ as a directed
limit of finite type $\mathbf{Z}$-algebras. By
Limits, Lemma \href{https://stacks.math.columbia.edu/tag/01ZM}{lemma descend finite presentation (Tag 01ZM)}
we can find an $i$ and a morphism $f_i : X_i \to \Spec(A_i)$ of
finite presentation, endowed with a section $s_i : \Spec(A_i) \to X_i$
whose base change to $Y$ recovers $f$ and the section $s$.
After possibly increasing $i$ we may also assume that $f_i$ is flat, see
Limits, Lemma \href{https://stacks.math.columbia.edu/tag/04AI}{lemma descend flat finite presentation (Tag 04AI)}.
Let $y_i$ be the image of $y$ in $Y_i$. Note that
$X_y = (X_{i, y_i}) \times_{y_i} y$. Hence $X_{i, y_i}$ is geometrically
reduced, see
Varieties, Lemma \href{https://stacks.math.columbia.edu/tag/0384}{lemma geometrically reduced upstairs (Tag 0384)}.
By
Lemma \href{https://stacks.math.columbia.edu/tag/055M}{lemma base change connected along section (Tag 055M)}
it suffices to prove the lemma for the system $f_i, s_i, y_i \in Y_i$.
Thus we reduce to the case where $Y$ is the spectrum of a Noetherian ring.

\medskip\noindent
Assume $Y$ is the spectrum of a Noetherian ring.
Since $f$ is of finite presentation,
i.e., of finite type, we see that $X$ is a Noetherian scheme too, see
Morphisms, Lemma \href{https://stacks.math.columbia.edu/tag/01T6}{lemma finite type noetherian (Tag 01T6)}.
Let $x \in X^0$ be a point lying over $y$. By
Topology, Lemma \href{https://stacks.math.columbia.edu/tag/0540}{lemma constructible neighbourhood Noetherian (Tag 0540)}
it suffices to prove that for any irreducible closed $Z \subset X$
passing through $x$ the intersection $X^0 \cap Z$ is dense in $Z$.
In particular it suffices to prove that the generic point $x' \in Z$
is in $X^0$. By
Properties, Lemma \href{https://stacks.math.columbia.edu/tag/054F}{lemma locally Noetherian specialization dvr (Tag 054F)}
we can find a discrete valuation ring $R$ and a morphism
$\Spec(R) \to X$ which maps the special point to $x$ and
the generic point to $x'$. We are going to think of $\Spec(R)$
as a scheme over $Y$ via the composition $\Spec(R) \to X \to Y$. By
Lemma \href{https://stacks.math.columbia.edu/tag/055M}{lemma base change connected along section (Tag 055M)}
we have that $(X_R)^0$ is the inverse image of $X^0$.
By construction we have a second section $t : \Spec(R) \to X_R$
(besides the base change $s_R$ of $s$)
of the structure morphism $X_R \to \Spec(R)$ such that
$t(\eta_R)$ is a point of $X_R$ which maps to $x'$ and
$t(0_R)$ is a point of $X_R$ which maps to $x$. Note that
$t(0_R)$ is in $(X_R)^0$ and that $t(\eta_R) \leadsto t(0_R)$.
Thus it suffices to prove that this implies that $t(\eta_R) \in (X_R)^0$.
Hence it suffices to prove the lemma in the case where $Y$
is the spectrum of a discrete valuation ring and $y$ its closed point.

\medskip\noindent
Assume $Y$ is the spectrum of a discrete valuation ring and $y$ is its closed
point. Our goal is to prove that $X^0$ is a neighbourhood of $X^0_y$.
Note that $X^0_y$ is open and closed in $X_y$ as $X_y$ has finitely
many irreducible components. Hence the complement $C = X_y \setminus X_y^0$
is closed in $X$. Thus $U = X \setminus C$ is an open neighbourhood of
$X^0_y$ and $U^0 = X^0$. Hence it suffices to prove the result for the
morphism $U \to Y$. In other words, we may assume that
$X_y$ is connected. Suppose that $X$ is disconnected, say
$X = X_1 \amalg \ldots \amalg X_n$ is a decomposition into connected
components. Then $s(Y)$ is completely contained in one of the $X_i$.
Say $s(Y) \subset X_1$. Then $X^0 \subset X_1$. Hence we may replace
$X$ by $X_1$ and assume that $X$ is connected. At this point
Lemma \href{https://stacks.math.columbia.edu/tag/055J}{lemma connected flat over dvr (Tag 055J)}
implies that $X_\eta$ is connected, i.e., $X^0 = X$ and we win.
\end{proof}
\end{document}
